\subsection{讨论}\label{sec:8.4}
\noindent\textbf{\autoref{sec:8.1}.} 超出 Hilbert 空间的背景，Banach 空间上的定量大数定律与型和余型（type and cotype）理论有关，见\cite{ref201}. 关于非正曲率空间上的重心，也可参见优秀论述\cite{ref293}，\autoref{exer:8.2}即取自其中.

\noindent\textbf{\autoref{sec:8.2}.} 本节材料取自\cite{ref4,ref199}.

\noindent\textbf{\autoref{sec:8.3}.} Wasserstein 重心的基本理论，包括存在性、对偶等，由\cite{ref3}发展，见\autoref{exer:8.3}. Wasserstein 重心的统计相合性由\cite{ref197}建立. \autoref{exer:8.4}中的方差不等式来自\cite{ref81}.

重心的计算也受到了大量关注. 对离散分布，\cite{ref9,ref10}的工作建立了固定维数下 Wasserstein 重心的多项式时间可解性，以及一般维数下的 NP 困难性；其他方法，例如通过神经网络参数化及连续优化方法的应用，见这些论文中的参考文献.

另一条与\autoref{sec:5}更密切相关的研究路线，是通过 Wasserstein 空间中的梯度方法计算重心\cite{ref13,ref14,ref28,ref53,ref81,ref189,ref329}. \autoref{exer:8.5}中的下降引理来自\cite{ref329}，它将\cite{ref14}的不动点方法解释为 Wasserstein 梯度下降.

高斯分布重心的研究早于一般情形，可以追溯到\cite{ref178,ref277}. 统计估计研究见\cite{ref188}，Wasserstein 梯度下降的非渐近计算保证则由\cite{ref13,ref81}给出.

与\autoref{sec:4}类似，可以在 Wasserstein 距离层面、重心目标函数层面，或者同时在两者中，为 Wasserstein 重心加入熵正则化，见\cite{ref37,ref65,ref82,ref187,ref207,ref312}.
